\section{Explicit affine block counts}
\label{sec:block-capacities}
Write $s_p=p^{v_p(s)}$, let $\operatorname{lpp}(t)$ denote the
largest prime-power divisor of $t>1$, and put
\[
 k(s)=\sum_qv_q(s)(q-1),\qquad
 \omega(s)=s\binom{k(s)}{\lfloor k(s)/2\rfloor}2^{-k(s)}.
\]

Use Lemma~\ref{lem:refined-induced-capacity} for the induced module
and its dual. In the notation of that lemma, write
\begin{equation}\label{eq:refined-H}
 H(s,p,a;T)=\min\left\{\bar H,
             \left\lceil\frac{4as}{\sqrt{\log s}}\right\rceil\right\},
\end{equation}
where $\bar H$ is the minimum of all its available integer bounds.
The soluble-transitive entries are always available at prime-power
block counts.

Always use \eqref{eq:refined-H} in
Theorem~\ref{thm:affine-chief}. Its upper cutoff retains the uniform
$\beta_U=o(w^2)$ and the semisimple error estimate. Thus a uniform
positive margin
\begin{equation}\label{eq:block-margin}
 \frac{h(w)-s}{8}-\eta_U\geq\rho w
\end{equation}
is precisely an application of Theorem~\ref{thm:parametric}, with
the actual top $T$ as quotient. In the following arithmetic
$\gamma=\log(3)/3<17/32$, using $3^{32}<2^{51}$.

\begin{theorem}[Regular prime and symmetric-three components]
\label{thm:prime-s3-capacity}
For an actual nonbinary action with exact local $C_p$ on $p$ points,
\eqref{eq:block-margin} holds with $\rho=1/4096$ except possibly at
the following block counts:
\[
 \begin{array}{c|l}
 p\geq5&\varnothing\\
 p=3&2,6\\
 p=2&\{2^a:1\leq a\leq8\}\cup\{3\cdot2^a:0\leq a\leq10\}.
 \end{array}
\]
For exact natural $S_3$, it holds with $\rho=1/1024$ outside
$\{2,3,4,6,8,12,18\}$.
\end{theorem}
\begin{proof}
For $C_p$ the coefficient is $H\log(p)/p$.
If $p\geq5$, use $H\leq s/2$ and $\log(p)/p\leq17/32$.
Subtracting the desired larger margin $w/64$ leaves at least
$((7p-25)s-8)/64>0$.

For $p=3$, if $s$ has a prime divisor at least five, use $H\leq s/5$;
the margin exceeds $23s/160-1/8$. If $4\mid s$, use $H\leq s/4$;
the margin exceeds $15s/128$. These cases leave only $s=3^b$ and
$s=2\cdot3^b$. At $s=3$, $H\leq1$ gives margin $>3/32$.
For $s=3^b,b\geq2$, the soluble bound gives $H\leq3s/8$, hence
margin $>13s/256-1/8$. For $s=2\cdot3^b,b\geq2$, use
$H\leq s/\sqrt6<7s/16$, hence margin $>9s/512$.
All exceed $3s/4096$. Only two and six remain.

For $p=2$, the margin is $s/8-H/2$. If the odd part has a
prime-power divisor at least five, $H\leq s/5$ leaves $s/40$.
For $s=2^a,a\geq9$, the central-binomial ratio is at most
$63/256$, leaving $s/512$. For $s=3\cdot2^a,a\geq11$, use
$H\leq s\sqrt{2/33}<127s/512$, leaving more than $s/1024$.
This exhausts all counts outside the displayed binary list.

For $S_3$, write $H_p=H(s,p,1;T)$. Its coefficient is
$H_2/2+\gamma H_3$ and its reserve is $s/4-(s\bmod2)/8$.
If $s$ has a prime divisor at least five, $H_2,H_3\leq s/5$
leave more than $7s/160-1/8\geq3s/160$.
For $s=2^a,a\geq4$, use $H_2\leq3s/8$, $H_3\leq1$;
the gap is greater than $s/16-17/32>3s/1024$.
At $s=9$, the integer bounds $H_2\leq1$, $H_3\leq3$ leave
more than $1/32>27/1024$. For $s=3^b,b\geq3$, use
$H_3\leq5s/16$, $H_2\leq1$; the gap exceeds
$43s/512-5/8>3s/1024$.
Finally let $s=2^a3^b$, $a,b\geq1$. For $a\geq3$,
$\eta/s<1/6+17/256$; for $a=2,b\geq2$,
$\eta/s<1/18+17/128$; for $a=1,b\geq3$,
$H_2\leq2,H_3\leq s/3$ give $\eta/s<1/54+17/96$.
Each is below $1/4-3/1024$. The untreated mixed cases are
$6,12,18$, proving the exact stated sufficient list.
\end{proof}

\begin{theorem}[Four-point component capacities]
\label{thm:four-local-capacity}
For exact natural $A_4$, \eqref{eq:block-margin} holds with
$\rho=1/4096$ outside $\{2,3,4,6,8,12,16,24\}$.
For exact natural $S_4$ it holds with the same $\rho$ outside
\[
 \{2^a:1\leq a\leq8\}\cup
 \{3\cdot2^a:0\leq a\leq10\}\cup\{5,10,15,20,30,60\}.
\]
These statements allow all proper, zero and nonsplit block kernels.
\end{theorem}
\begin{proof}
Put $H_{pa}=H(s,p,a;T)$. The respective coefficients are
$\eta_A=H_{22}/2+\gamma H_{31}$ and
$\eta_S=\eta_A+H_{21}/2$, and the reserve is $3s/8$.

For $A_4$, a prime divisor at least five gives
$\eta_A/s<(1+17/32)/5=49/160$.
Otherwise $s=2^a3^b$. For pure binary powers $a\geq5$,
$H_{22}\leq5s/8,H_{31}\leq1$ leave more than $s/16-17/32$.
For pure ternary powers $b\geq2$,
$H_{22}\leq2,H_{31}\leq3s/8$ leave more than $45s/256-1$.
For $a,b\geq1$ with $a,b\geq2$ use
$\eta_A/s<1/9+17/128$; for $a=1,b\geq2$ use
$H_{31}\leq s/\sqrt6<5s/12$, giving
$\eta_A/s<1/9+85/384$. Finally $b=1,a\geq4$ gives
$\eta_A/s<1/3+17/512$, whose gap is $13/1536$.
All displayed gaps are at least $s/1024$, and their complement
is exactly the asserted eight counts.

For $S_4$, a prime divisor at least seven gives
$\eta_S/s<65/224$; a factor $25$ gives the still smaller bound
$65/800$. Write therefore $s=2^a3^b5^c$, $c\in\{0,1\}$.
When $c=1$, the prime-to-characteristic estimates give
\[
 \eta_S/s\leq\frac{3/2}{\max(3^b,5)}+
                         \frac{\gamma}{\max(2^a,5)}.
\]
For $a\geq3$ this is less than $3/10+17/256$, leaving
$11/1280$; for $b\geq2$ it is less than $1/6+17/160$.
The six complementary counts are $5,10,15,20,30,60$.
When $c=0$, pure binary powers $a\geq9$ give
$\eta_S\leq189s/512+\gamma$, leaving more than
$3s/512-17/32\geq s/1024$. Pure ternary powers $b\geq2$
leave more than $45s/256-3/2\geq s/1024$.
For mixed counts, $b\geq3$ gives
$\eta_S/s<1/18+17/64$, and $a,b\geq2$ gives
$\eta_S/s<1/6+17/128$. At the remaining $s=18$, the integer
bounds $(H_{22},H_{21},H_{31})\leq(4,2,7)$ leave more than
$1/32>18/1024$. Finally $b=1,a\geq11$ gives
\[
 H_{22}\leq2s\sqrt{2/33}<79s/160,\quad
 H_{21}<79s/320,\quad H_{31}\leq3;
\]
the gap exceeds $3s/640-51/32\geq s/1024$ for $s\geq6144$.
The stated exceptional set is exactly the complement of these cases.
\end{proof}

\begin{theorem}[The four remaining affine local degrees]
\label{thm:small-affine-capacity}
For $r\in\{5,8,9,16\}$, every action with exact primitive affine
local degree $r$ satisfies \eqref{eq:block-margin} with
$\rho=1/4096$, except possibly at
\[
 \begin{array}{c|l}r&s\\\hline
 5&2,3,4\\8&2,4,6\\9&2,3,4,6,12\\16&2,4.
 \end{array}
\]
\end{theorem}
\begin{proof}
Let $H(s,p,a)$ be the minimum in Lemma~\ref{lem:refined-induced-capacity},
using the soluble entries only for prime-power $s$ and omitting the
logarithmic cutoff. It has the form $\lfloor ax(s,p)\rfloor$,
so $\sum_jH(s,p,a_j)\leq H(s,p,\sum_ja_j)$. Floors are taken
after multiplying by the dimension. Put $t_2=1/2,t_3=17/32$,
$t_5=7/15,t_7=3/7$, which upper-bound $\log p/p$.
For brevity write $H_{pa}=H(s,p,a)$. The following upper bounds
retain the true translation dimension:
\begin{align*}
 E_5&=t_5H_{51}+\tfrac12H_{22},\\
 E_8&=\tfrac12H_{23}+t_3H_{31}+t_7H_{71},\\
 E_9&=\tfrac12H_{24}+t_3(H_{32}+H_{31}),\\
 E_{16}&=\tfrac12H_{24}+
   \max_{\boldsymbol n\in V}\sum_{p=2,3,5,7}t_pH_{p n_p},
\end{align*}
where
$V=\{0\leq\boldsymbol n\leq(6,2,1,1):\sum n_p\leq5\}$.
Indeed the stabilizers at degrees five and nine divide $C_4$ and
$\mathrm{GL}_2(3)$, respectively. At degree sixteen the order
$|\mathrm{GL}_4(2)|=2^6 3^2\cdot5\cdot7$ gives the coordinate
bounds, and \cite[Theorem~1.3]{GlasbyPraegerRosaVerret2018}
gives at most five stabilizer composition factors after subtracting
the four translation factors.

For degree eight, a proper subgroup of the simple group
$\mathrm{GL}_3(2)$ has index at least seven, since an index below
seven would embed the group in $S_6$. Its order divides 168 and is
at most 24, and therefore has at most two prime divisors and is
soluble. An irreducible soluble stabilizer has trivial normal
$2$-subgroup. A minimal normal subgroup is then $C_3$ or $C_7$;
$C_3$ has a nonzero proper fixed space on $\F_2^3$, contradicting
irreducibility. A normal $C_7$ puts the stabilizer between $C_7$
and its order-21 Singer normalizer. Thus the only possibilities are
$C_7,C_7:C_3,\mathrm{GL}_3(2)$, proving the envelope $E_8$.
Nonabelian local factors in these arguments stay in $\Psi_U$.

Here is an exhaustive arithmetic verification for all $s$, without
a group census. Independent upper bounds on weighted dimensions and
a binary exponent cutoff are
\[
 \begin{array}{c|rrrr|r}
 r&A_2&A_3&A_5&A_7&a_0\\\hline
 5&1&0&7/15&0&5\\
 8&3/2&17/32&0&3/7&3\\
 9&2&51/32&0&0&4\\
 16&9/2&17/16&7/15&3/7&5.
 \end{array}
\]
The reserve per $s$ is at least $B_r=(r-1)/8-1/16$ for odd $r$,
and $B_r=(r-1)/8$ for even $r$. If $5\mid s$, the capacity per
dimension at $p=5$ is less than $5s/12$, and at the other primes
at most $s/5$. If $7\mid s$, use $s/3$ at seven and $s/7$
elsewhere. A prime divisor at least eleven gives $s/11$ everywhere.
If $9\mid s$, use $5s/12$ at three and $s/9$ elsewhere.
If $2^{a_0}\mid s$, use $3s/8$ at two for $r=5,16$, $s/2$
for $r=8$, and $5s/12$ for $r=9$; other primes get $s/2^{a_0}$.
Substitution gives these lower gaps per $s$, respectively:
\[
\begin{array}{c|rrrrr}
r&5\mid s&7\mid s&q\geq11&9\mid s&2^{a_0}\mid s\\\hline
5&31/720&383/1680&73/240&593/2160&23/480\\
8&429/1120&99/224&1605/2464&1181/2688&9/1792\\
9&7/32&95/224&215/352&59/1152&7/1536\\
16&2431/5040&209/240&23801/18480&50369/60480&6791/53760.
\end{array}
\]
Every entry exceeds $r/4096$. The complement is exactly
$s=2^a3^b$, $0\leq a<a_0$, $b\in\{0,1\}$, $s\geq2$.
The accepted cells there, with exact gaps $(h(rs)-s)/8-E_r$, are
\[
\begin{array}{c|l}
5&(6,1/15),(8,8/15),(12,3/5),(16,23/15),(24,13/5),(48,33/5)\\
8&(3,37/224),(12,363/224)\\
9&(8,13/32),(24,103/32)\\
16&(3,17/32),(6,17/16),(8,3/2),(12,9/2),(16,3),(24,9),(48,18).
\end{array}
\]
Direct substitution into the displayed integer formula proves each
entry; each exceeds $rs/4096$. The thirteen remaining cells are
exactly the exceptions in the theorem. Taking the logarithmic
cutoff again restores the uniform error estimates and cannot worsen
any margin.
\end{proof}

\begin{theorem}[Zero ternary kernel at local degree four]
\label{thm:zero-ternary-capacity}
Let $U$ have exact local $A_4$ or $S_4$ on $s\geq2$ blocks, put
$M=(\F_2^2)^s$, $E=U\cap M$, and view $R=U/E$ inside
$R_0\wr T$, where $R_0=C_3$ or $S_3$. If $R\cap C_3^s=1$,
then $R=T$, counts two and three are impossible, and every count
outside $\{4,8,16\}$ satisfies \eqref{eq:block-margin} with
$\rho=1/96$.
\end{theorem}
\begin{proof}
For local $A_4$, the linear block kernel is already $R\cap C_3^s$.
For local $S_4$, it embeds in $C_2^s$. Each of its coordinate images
is a normal $2$-subgroup of the exact $S_3$, and is trivial. Thus
$R=T$ in both cases. At counts two and three its point stabilizer
cannot surject onto the required local $C_3$ or $S_3$.

Only the translation intersection layer remains. With
$H=H(s,2,2;T)$ its entire fibre is bounded by
\[
 Z_J(U)\leq2^{Hb/2+4sH+2s(g(4s)+1)}Z_J(T).
\]
If $s$ is not a power of two, $H\leq2s/3$ leaves gap
$3s/8-H/2\geq s/24=w/96$. If $s=2^a,a\geq5$,
$H\leq2\binom a{\lfloor a/2\rfloor}\leq5s/8$ leaves $s/16$.
The other powers are two, four, eight and sixteen, of which two is
impossible. The cutoff in \eqref{eq:refined-H} supplies the uniform
quadratic error, so Theorem~\ref{thm:parametric} applies.
\end{proof}
